La gestión de datos maestros de materiales es un proceso crítico para cualquier organización que opera sobre SAP, ya que la calidad de este catálogo incide directamente en procesos como la ejecución de órdenes de compra, el cierre de inventarios y la trazabilidad de activos. El presente proyecto se desarrolla sobre el módulo de materiales en SAP de una unidad de negocio del sector energético que opera en varios países de Centroamérica y el Caribe, donde un equipo de gestores de datos maestros es responsable de validar que cada material nuevo no exista ya en el sistema, que esté correctamente categorizado y que cumpla con las especificaciones mínimas exigidas por la normativa ISO 8000.

En la práctica, este proceso presenta tres problemas recurrentes que se identificaron y cuantificaron mediante un análisis exploratorio sobre 39,571 materiales del catálogo de la unidad de negocio, distribuidos en 1,234 categorías:

\begin{enumerate}
	\item \textbf{Materiales duplicados.} Se identificaron 456 casos de duplicidad exacta en el catálogo, y un riesgo potencial de duplicidad parcial (por similitud textual) que alcanza al 94\% de los registros, producto de descripciones libres, abreviaciones inconsistentes entre gestores y el límite de 40 caracteres que impone SAP para el campo de descripción corta.
	\item \textbf{Categorización incorrecta.} 246 materiales presentan una clasificación inadecuada respecto a su naturaleza real, y 94 repuestos específicamente fueron creados bajo la categoría de ``suministros'', lo que distorsiona los reportes de consumo y dificulta la planificación de compras por tipo de material.
	\item \textbf{Tiempo de gestión elevado.} Con un volumen promedio de 4.4 solicitudes de creación individual por día hábil (medido entre el 15 de febrero y el 17 de junio de 2026, período previo a la implementación), cada gestión permanece en promedio 3.3 horas dentro del rol del gestor —tiempo que incluye cola de trabajo, validaciones y los distintos pasos del proceso—. De ese tiempo, la búsqueda manual en SAP para descartar duplicados y verificar especificaciones —el paso que GAMMA automatiza directamente— consume entre 5 y 10 minutos por solicitud.
\end{enumerate}

Estos hallazgos evidencian que el problema no es la falta de un proceso de gestión de datos maestros —la organización ya cuenta con uno—, sino la ausencia de herramientas analíticas que apoyen al gestor en el momento de la decisión: al crear un material, no existe hoy ningún mecanismo automatizado que le indique si ese material ya existe, ni qué categoría le corresponde con base en el histórico del catálogo.

Para atender esta necesidad se desarrolló \textbf{GAMMA — Gobierno Automatizado del Maestro de Materiales}, un asistente inteligente que integra tres módulos complementarios: normalización de descripciones mediante un modelo de lenguaje, detección de duplicados por similitud difusa sobre el catálogo, y clasificación automática de categoría mediante un modelo de aprendizaje supervisado entrenado sobre el histórico del catálogo. El proyecto abarca el ciclo de vida completo de un proyecto de analítica: desde el análisis exploratorio de los datos de origen, pasando por la integración y modelado, hasta el despliegue de una herramienta funcional validada por el propio equipo de gestores.

El presente documento describe la metodología seguida en cada una de estas etapas —análisis de fuentes de información, integración de datos, explotación de la información, desarrollo de modelos analíticos, visualización de resultados y arquitectura de puesta en producción— así como los resultados obtenidos y su impacto en la operación, tanto a nivel de eficiencia operativa como de retorno financiero.